\section{Conclusion}\label{sec:conclusion}

Evaluative bibliometrics and global university ranking systems frequently conflate national research assessment corpora with address-indexed bibliographic databases. Linking the comprehensive submission dataset of Australia's ERA 2018 assessment against OpenAlex refutes this equivalence. Across 2011--2016, 46.1\% (343,313 works) of all publications bearing an Australian university author affiliation in OpenAlex were absent from the national assessment submission. This omitted corpus includes 222,428 standard peer-reviewed outputs: 169,636 journal articles, 29,284 book chapters, 20,032 conference papers, and 3,476 books.

These omitted works are intrinsically integrated into the national scientific literature: ERA submissions direct 109,443 citations into 45,822 unsubmitted Australian works, encompassing foundational algorithmic and methodological contributions recording tens of thousands of global citations. Furthermore, among submitted works, the institution claiming the paper in ERA matches the published institutional address on the work in only 67.64\% of cases. A total of 49,962 submitted works (13.76\%) carry zero Australian institutional affiliations on the published document. This discordance follows a strict monotonic temporal gradient across the six-year evaluation window, with pure foreign attribution expanding from 8.40\% in 2016 to 19.44\% in 2011.

These disparities are not artifacts of bibliographic indexing failure or imperfect string matching. Grounded in Olensky's taxonomy of data accuracy in bibliometric sources \cite{olenskyDataAccuracyBibliometric2015}, these patterns demonstrate the structural divergence between two incompatible bibliometric constructs: an \textbf{address-based production model} measuring research performed at physical institutional sites, and a \textbf{census-based workforce evaluation model} assessing the mobile intellectual output of personnel employed on a single statutory day. Evaluative exercises and citation impact studies must explicitly account for workforce portability, student attrition, and hospital/institute payroll boundaries rather than treating assessment submissions as complete proxies for institutional research output.
